\section{The Solidity and Vyper Backends}
\label{sec:solidity}

The EVM backends translate the EventGraph into a deployable smart
contract.  Both languages target the same compilation pass via a
shared EVM intermediate representation (\TT{evm.EVMIR}); only the
final pretty-printer differs.  This section walks through the
Solidity output and notes the few divergences for Vyper.  We are
careful to describe the actual generated code, not an idealised
version of it: the bail-out mechanism in particular is simpler than
the language-level story might suggest, and a few of its
consequences belong in the open-questions list rather than the
guarantees list.

\subsection{Threat model}
\label{sec:solidity:threats}

The compiled contract is meant to survive a transaction-ordering
adversary --- a builder or competing searcher who can choose which
pending transactions to include and in what order --- and to remain
liveness-preserving when a participant stops responding.  It is not
designed against:

\begin{itemize}\itemsep1pt
\item \emph{Censorship.}  An adversary who indefinitely excludes a
  participant's transactions is modeled as forcing a timeout, which
  marks the participant as \TT{bailed} and triggers the language-level
  quit handler.
\item \emph{Cross-contract interactions.}  The contract is a closed
  system; no calls into other contracts are emitted (the only
  outbound action is a direct ETH transfer to a registered
  participant address during withdrawal).
\item \emph{Gas-cost games.}  Gas is a real-world cost that affects
  which moves are rational, but the analytical backends ignore it.
  Bounded per-action gas envelopes (e.g.\ via
  GASTAP~\citep{Albert2020Gastap}) would be needed to bring gas into
  the strategic model; this is planned but not implemented.
\item \emph{Cryptographic primitives.}  The contract relies on
  keccak-256 being binding and hiding.  Domain separation and
  per-deployment binding are in the hash composition
  (\S\ref{sec:solidity:commit-reveal}), but the strength of the
  primitive itself is taken as given.
\end{itemize}

The remainder of this section describes the construction; the
``honestly imperfect'' parts of the design (especially the single
global timeout window) are flagged inline.

\subsection{Storage layout}

The contract's storage is derived directly from the EventGraph.
There are five families of slots.

\paragraph{Per-node completion.}  A two-level mapping
\TT{actionDone[role][actionId]} records, for each role and each
node, whether the node has fired.  A symmetric mapping
\TT{actionTimestamp[role][actionId]} stores the block timestamp at
which it fired.  Node ids are integers assigned by a deterministic
topological linearisation of the EventGraph; one
\TT{uint256 constant ACTION\_\textit{Role}\_\textit{n}} is emitted
per node.  The integer \TT{\textit{n}} is the node's source position
(making the constant name human-readable), while the assigned
constant \emph{value} is its topological index.

\paragraph{Role registry.}  A mapping
\TT{roles[address] => Role} associates EOAs with the roles they
have joined, and a per-role \TT{address\_\textit{Role}} slot
stores the registered address.  A small \TT{enum Role} carries one
constructor per role plus a \TT{None} sentinel.

\paragraph{Game state.}  For each \emph{written field} in the
EventGraph the contract stores the field's value.  A field with
\PUBLIC visibility uses a slot of the field's native type
(\TT{int256}, \TT{bool}, etc.); a field with \COMMIT visibility
uses a separate \TT{bytes32} slot for the commitment hash, with the
revealed value filled in by the matching reveal action.  Each
value slot is shadowed by a \TT{done\_\textit{slot}: bool} flag.

\paragraph{Per-role bookkeeping.}  Three booleans per role:
\TT{done\_\textit{Role}} (the role has joined),
\TT{claimed\_\textit{Role}} (the role has claimed its terminal
payout), and a private \TT{bailed[\textit{Role}]} (the role has been
marked inactive after a timeout).

\paragraph{Game-wide constants.}  A constant \TT{TIMEOUT} sets the
inactivity window, in seconds (the default is 24 hours).  A
constant \TT{COMMIT\_TAG} (a hashed domain-separation tag) and the
contract's deployed address are folded into every commitment hash.

\subsection{Modifiers and the depends pattern}
\label{sec:solidity:depends}

The contract has no global step counter.  Each per-node function is
guarded by a \TT{depends} modifier per predecessor.  The modifier
itself does two things: it advances the contract's
inactivity-timeout state (potentially marking some role
\TT{bailed}), and it conditionally requires the predecessor to have
completed.  Read top to bottom:

\begin{lstlisting}[style=sol]
modifier depends(Role role, uint256 actionId) {
    _check_timestamp(role);
    if (!bailed[role]) {
        require(actionDone[role][actionId], "dependency not satisfied");
    }
    _;
}
\end{lstlisting}

The body invokes \TT{\_check\_timestamp(\textit{role})}, which marks
\TT{\textit{role}} as \TT{bailed} iff the current
\TT{block.timestamp} exceeds a shared global high-water mark
\TT{lastTs} by more than \TT{TIMEOUT}; otherwise it leaves the
state unchanged.  If \TT{\textit{role}} is not bailed, the
predecessor's completion flag is required; if it is bailed, the
\TT{require} is skipped and the dependent action proceeds.  Two
additional modifiers complete the gating: \TT{action(role, id)}
asserts that the action has not yet fired, marks it done, runs the
body, and updates \TT{actionTimestamp} and \TT{lastTs};
\TT{by(role)} asserts that the calling EOA holds the given role
and is not itself bailed.

A typical per-node function looks like:

\begin{lstlisting}[style=sol]
function move_Guest_3(int256 _d)
    external
    by(Role.Guest)
    depends(Role.Host, ACTION_Host_0)
    depends(Role.Guest, ACTION_Guest_1)
    depends(Role.Host, ACTION_Host_2)
    action(Role.Guest, ACTION_Guest_3)
{
    require(_d >= 0 && _d <= 2, "domain");
    Guest_d = _d;
    done_Guest_d = true;
}
\end{lstlisting}

Every emitted action is named \TT{move\_\textit{Role}\_\textit{n}}
where \textit{n} is the source position.  The function's dependency
list is exactly the predecessor set of the node in the EventGraph
(after the commit-reveal expansion).  Concurrent nodes share the
prefix of the chain up to the divergence point and are otherwise
free to be called in any order.

\subsection{Commits and reveals}
\label{sec:solidity:commit-reveal}

Every COMMIT-tagged write becomes a function that takes a
\TT{bytes32} commitment and stores it; every REVEAL-tagged write
takes the cleartext value plus a salt, recomputes the commitment,
and asserts equality with the stored hash before performing the
public write.  The commitment hash is composed as

\begin{equation*}
  \mathit{keccak256}\bigl(
    \mathrm{abi.encode}(
      \textsc{COMMIT\_TAG},\
      \mathrm{address(this)},\
      \mathrm{role},\
      \mathrm{actor},\
      \mathit{keccak256}(\mathrm{payload})
    )
  \bigr).
\end{equation*}

The pieces serve distinct roles.  The \texttt{COMMIT\_TAG} provides
cross-protocol domain separation.  \TT{address(this)} ties the
commitment to this particular deployment, preventing replay across
copies of the same protocol.  The \TT{role} identifier defends
against role-swapping replays within a single deployment; the
\TT{actor} address binds the commitment to the EOA that posted it.
The payload is pre-hashed before being folded in, so the outer
hash works over a fixed-size input regardless of the parameter
list.

Vyper emits the analogous construction using the language's
\TT{assert} idiom and slightly different storage syntax.  The
generated bytecode is comparable.

\subsection{The single-window timeout and its consequences}
\label{sec:solidity:bail}

The timeout mechanism is simpler than the language-level story
might suggest, and it is worth describing exactly.

The contract maintains a single global \TT{lastTs} slot, updated to
\TT{block.timestamp} on every successful action.  Whenever a
modifier is entered, \TT{\_check\_timestamp(role)} is called: if
\TT{block.timestamp > lastTs + TIMEOUT}, the named role is marked
\TT{bailed} \emph{and} \TT{lastTs} is reset to the current
timestamp.  In effect, the contract has one sliding inactivity
window shared across all participants: if no action has happened
anywhere for \TT{TIMEOUT} seconds, the next caller can mark
\emph{some} role as bailed and the window restarts.

This has two non-obvious consequences worth flagging:

\begin{itemize}\itemsep1pt
\item \emph{The role that gets marked is the role named in the
  caller's modifier.}  \TT{depends(Role.Host, \ldots)} on a window
  expiry marks \TT{Host} as bailed; \TT{by(Role.Guest)} on a window
  expiry marks \TT{Guest}.  This is approximately right --- the
  function call points at a role expected to have acted --- but it
  is not the same thing as ``the unique unresponsive role''; the
  attribution is structural rather than analytical.
\item \emph{Bail is one-way.}  Once \TT{bailed[\textit{R}] = true},
  the \TT{by(R)} modifier rejects any subsequent action from
  \TT{R}.  A role that goes offline briefly and reconnects after
  the window expires has lost the right to play.  The
  language-level treatment of \TT{Quit} as a definitive move is the
  source of this design; it is operationally simple, but it does
  mean that a fair-playing role with a flaky connection is at risk
  in long protocols.
\end{itemize}

When a dependent action runs against a bailed predecessor, the
\TT{depends} modifier skips its \TT{require}, but it does \emph{not}
fabricate writes for the missing fields.  Their value slots and
their \TT{done\_} flags remain in their default state (zero, and
\TT{false} respectively).  Downstream code that needs to
distinguish ``never written'' from ``written to the default value''
relies on the language's \TT{opt T} typing: a field declared with
an \TT{or null} handler is typed \TT{opt T}, and the \TT{withdraw}
clause is required to inspect its \TT{done\_} flag.

The fair-play guarantee one would like --- ``every role can drive
the contract from any reachable configuration to a terminal
allocation'' --- is plausible for protocols whose \TT{withdraw}
clause is total on the public state (including \TT{done\_} flags),
but it is not proved.  Stating and proving it under the explicit
execution model is the central refinement-theorem target of the
ongoing VegasCore work; today, the test suite checks that the
\TT{withdraw} clause type-checks against the optional fields it
must handle, and that hand-played example traces reach a terminal
allocation, but no machine-checked proof is shipped.

\subsection{Withdraw and payout}

The terminal \TT{withdraw} clause is compiled into one entry
point per role.  Each role calls its own
\TT{withdraw\_\textit{Role}()} to receive their share, guarded by
a \TT{claimed\_\textit{Role}} flag to prevent double-payment.  The
payout expression is lowered to EVM IR by direct structural
translation; the body computes the expression against the public
state (including \TT{done\_} flag inspection for optional fields)
and transfers the resulting amount to the role's stored address
using a low-level \TT{call} with the computed value.

Conservation is checked statically against the deposit total
during compilation (symbolic checks where possible; exhaustive
enumeration for the finite-type case).  The contract therefore
does not need to verify at runtime that all role payouts sum to
the pool.

\subsection{Things the backend deliberately leaves alone}

Several real concerns are not addressed by the Solidity emission
because they are scoped out of the language.  Gas costs as a
strategic factor; mempool-level censorship beyond what the bail
mechanism handles; cross-contract interactions and oracle reads;
admin keys, upgrade proxies, and pausable patterns.  Each is a
deliberate restriction in the same vein as the language's
bounded-types-only stance: drawing the boundary tightly is what
makes the analytical backends' outputs faithful to the deployed
code, and what makes ``the same source for two views'' a defensible
claim.
